\documentclass[12pt]{article}
\usepackage[utf8]{inputenc}
\usepackage[spanish]{babel}
\usepackage{graphicx}
\usepackage{float}
\usepackage{caption}
\usepackage{csquotes}
\usepackage{booktabs}
\usepackage{hyperref}
\usepackage[backend=biber,style=apa]{biblatex}
\addbibresource{../../referencias/referencias.bib}
\usepackage{geometry}
\usepackage{url}
\geometry{margin=2.5cm}



\begin{document}
\begin{titlepage}
\centering
\vspace*{\fill}

{\Huge\bfseries Bitácora 3\par}
\vspace{1.5cm}

{\Large\textbf{Integrantes:}\par}
Debbie Con Ortega (C32250)\\
Ashly Garro Villanueva (C33198)\\
Emily Sanchéz Mancía (C27260)\\
Alessandro Umaña Vega (C37963)

\vspace{1cm}

{\Large\textbf{Profesor:} Maikol Solís\par}
\vspace{1.5cm}

{\Large I Ciclo 2026\par}

\vspace*{\fill}
\end{titlepage}
\tableofcontents
\newpage

\section{Rastro de decisiones}
\section{Registro de uso de IA}
\section{Log de contribuciones}
\section{Respuesta a retroalimentación} 
\section{Audiencia del proyecto} 





\section{Parte de planificación}
\subsection{Fichas de literatura}
\textbf{Ficha 1:}

\textbf{Nivel descriptivo}
\begin{itemize}
    \item \textbf{Título:} Traffic Accidents Analysis on Dry and Wet Road Bends Surfaces in Greater Manchester-UK. 
    
    \item \textbf{Autores:} Ismael Karzan S, Razzaq Noorance A. 
    
    \item \textbf{Año:} 2017 
    
    \item \textbf{Nombre del tema:} Accidentes de tránsito en curvas de Greater Manchester en función de condiciones 
    climáticas y factores de entorno como iluminación, velocidad y comportamiento humano. 
    
    \item \textbf{Clasificación:} Temático. El artículo hace uso de herramientas estadísticas (como la prueba Chi-Cuadrado) y 
    análisis descriptivo para estudiar la relación entre distintas variables, lo cual es directamente aplicable al enfoque 
    de nuestro proyecto. A pesar de esto, se consideran más importantes sus conclusiones para identificar los factores que 
    afectan en mayor magnitud la frecuencia de accidentes de tránsito, por lo que se clasifica como temático.
    
    \item \textbf{Resumen en una oración:} Analiza cómo condiciones de carretera afectan accidentes bajo distintas condiciones 
    climáticas y de entorno. 
    
    \item \textbf{Argumento central:} Las condiciones de la superficie vial (seca o húmeda), junto a otros factores como la 
    iluminación, velocidad y las decisiones del conductor, influyen significativamente en la ocurrencia de accidentes, y estas 
    relaciones pueden evaluarse mediante análisis descriptivo y métodos estadísticos. 
    
    \item \textbf{Problemas con el argumento o el tema:} Se limita a una región específica (Greater Manchester), lo que reduce 
    la generalización de las conclusiones. Por otra parte, considera condiciones climáticas limitadas, dejando por fuera 
    situaciones de alto riesgo como la nevisca. Además, algunas relaciones, como la superficie vial y la severidad de los 
    accidentes, no resultan estadísticamente significativas lo que limita conclusiones fuertes.  \\
\end{itemize}

\textbf{Nivel analítico}

\begin{itemize}
    \item \textbf{Conexión con mi proyecto:} El artículo es sumamente relevante ya que analiza, mediante pruebas estadísticas 
    aplicables para el presente proyecto, variables similares a las consideradas para esta investigación, como condiciones 
    climáticas e iluminación, usando datos reales del Reino Unido (STATS19). 
    
	\item \textbf{Lo que NO dice el autor:} A diferencia del artículo, este estudio busca explicar cómo varía la cantidad
    total de accidentes en el Reino Unido (no solo en curvas de Greater Manchester), utilizando la base de datos UK Road 
    Safety: Traffic Accidents and Vehicles. Además, no solo se describirán factores asociados (como clima o condición de la 
    carretera), sino que se buscará identificar situaciones de riesgo recurrentes para proponer medidas de prevención. 
    
	\item \textbf{Contraste con otra fuente:} El estudio “Accident risk of road and weather conditions on different road 
    types” por Malin et al. analiza el riesgo relativo de accidentes considerando múltiples condiciones climáticas (nieve, 
    hielo, visibilidad, temperatura) y distintos tipos de carretera, incorporando además la exposición real al riesgo mediante 
    probabilidad de Palm. En contraste, el estudio de Greater Manchester se enfoca en identificar asociaciones entre variables
    categóricas (superficie seca y húmeda, iluminación, velocidad) mediante pruebas Chi-Cuadrado en un contexto local. 
    Sin embargo, ambos coinciden en que las condiciones de la vía y el entorno influyen significativamente en la ocurrencia 
    de accidentes, y en que el análisis estadístico de estas variables permite identificar factores de riesgo relevantes, 
    aunque con distintos niveles de profundidad.  
    
	\item \textbf{Evaluación de aplicabilidad (1-5):} 4/5. La metodología del artículo es compatible con la presente investigación 
    debido a que utiliza pruebas de independencia sobre variables categóricas, que coinciden con la estructura de la base 
    de datos utilizada para este. Sin embargo, no cumple todos los puntos de aplicabilidad debido a que el 
    enfoque del autor es local y específico a curvas en Greater Manchester mientras que el actual estudio abarca todo el 
    Reino Unido y un conjunto de factores más amplio, lo que exige adaptar y ampliar el diseño analítico del artículo.
    
	\item \textbf{Pregunta que le haría al autor:} Dado que su análisis se centra en la relación de la superficie vial en 
    curvas de Greater Manchester tomando en cuenta el factor humano para estudiar las causas de estos, ¿cómo ajustaría su 
    metodología si se quisiera estudiar solamente factores ambientales e infraestructurales, excluyendo el comportamiento 
    del conductor, y qué variables consideraría más relevantes para poder aplicar su enfoque a una base más amplia como la 
    utilizada en este estudio?
    
	\item \textbf{Resumen argumentativo:} Los autores buscan los factores que influyen en la ocurrencia de accidentes de 
    tránsito en Greater Manchester, particularmente en curvas, analizando asociaciones entre variables como la superficie 
    de la carretera, iluminación, velocidad y severidad del accidente, utilizando pruebas de independencia como la Chi-Cuadrado.
    La solución propuesta resulta parcialmente aplicable a el contexto de este estudio, ya que el enfoque metodológico basado en el
    análisis de variables categóricas y la identificación de asociaciones entre factores ambientales es coherente con los 
    objetivos del proyecto. Sin embargo, presenta limitaciones importantes, dado que el estudio se restringe a curvas 
    de Greater Manchester, y considera el comportamiento humano como un factor central, mientras que en esta investigación
    se centra en variables climáticas y del entorno, analizando accidentes en todo el Reino Unido. En este trabajo 
    se retomará el enfoque de análisis asociativo entre variables, por ejemplo, entre condiciones climáticas y cantidad 
    de accidentes, así como la identificación de factores de riesgo, sin embargo, se aplicará a un conjunto de datos 
    más amplio y sin atribuir la causalidad principal al error del conductor.\\
\end{itemize} 

\textbf{Ficha 2:}

\textbf{Nivel descriptivo}
\begin{itemize}
    \item \textbf{Título:} The effects of drivers' speed on the frequency of road accidents. 
    
    \item \textbf{Autores:} M. C. Taylor, D. A. Lynam y A. Baruya. 
    
    \item \textbf{Año:} 2000. 
    
    \item \textbf{Nombre del tema:} Relación cuantitativa entre la velocidad de los conductores y la frecuencia de accidentes en diferentes tipos de carreteras. 

    \item \textbf{Clasificación:} Teórica. Se elige esta clasificación porque que sustenta el argumento de que la velocidad es un factor determinante en la variación de la frecuencia de accidente. 

    \item \textbf{Resumen en una oración:} Investigación sobre cómo la velocidad media y su variación afectan la frecuencia de accidentes viales. 

    \item \textbf{Argumento central:} El documento busca demostrar que la velocidad a la que circulan los vehículos está relacionada no solo con la gravedad de los accidentes de tránsito, sino también con su frecuencia. El estudio se enfoca principalmente en la velocidad promedio de los vehículos en una vía determinada, en lugar de la velocidad de un conductor individual. A partir de este enfoque, el estudio analiza información proveniente de carreteras urbanas y rurales e investigaciones sobre el comportamiento de los conductores con el objetivo de identificar patrones claros entre velocidad y accidentes. Los resultados muestran que existe una relación positiva entre aumentos en la velocidad promedio y aumentos en el número de accidentes. Esto se explica porque la velocidad influye en factores clave del proceso de conducción, como el tiempo de reacción del conductor, la distancia de frenado y la capacidad para evitar choques. También, el documento señala que incluso pequeños cambios en la velocidad promedio pueden producir cambios relativamente grandes en la frecuencia de accidentes y que el efecto de estas variaciones depende también del tipo de carretera y de las condiciones de circulación. Por esta razón, los autores concluyen que la regulación  de la velocidad es un elemento fundamental en las políticas de seguridad vial, aunque también reconocen que otros factores, como las condiciones de la carretera, el clima, el diseño de las intersecciones, pueden influir en el riesgo de accidentes. 
 
    \item \textbf{Problemas con el argumento o el tema:} Los autores mencionan la complejidad de analizar la relación entre la velocidad y la frecuencia de los accidentes, pues reconocen el efecto que otras variables presentan sobre los mismos. Por este motivo, aislar el efecto de la velocidad no es sencillo. Por otro lado, se menciona la dificultad de conocer con precisión la velocidad real del vehículo involucrado en un accidente, por lo que, en su lugar, se toma la velocidad promedio de los vehículos en la vía.  \\
    \end{itemize}

\textbf{Nivel analítico}

\begin{itemize}
  \item \textbf{Conexión con mi proyecto:} El estudio se relaciona directamente con el proyecto que se desea realizar, pues no solo también se hizo en el Reino Unido, si no también provee evidencia empírica de que la velocidad sí afecta la probabilidad de tener un accidente. En particular, el análisis que se desea realizar tiene como variables la velocidad límite de la carretera donde se dio el accidente, el tipo de carretera, si habian intersecciónes, el área, etc. Por ende, la lectura índica que teóricamente que los accidentes ocurren con mayor frecuencia en carreteras urbanas, especialmente en vías con mucho tráfico, presencia de peatones y muchas intersecciones, porque hay más interacciones entre los diferentes usuarios de la vía. En cambio, en carreteras rurales los accidentes suelen ocurrir con menor frecuencia, aunque cuando suceden pueden ser más graves. Por lo que se esperaría que los resultados del análisis de datos sean congruentes con el artículo.  
 
	\item \textbf{Lo que NO dice el autor:} El estudio no analiza factores climáticos específicos, como lluvia, ni cómo estas condiciones pueden afectar la frecuencia de accidentes. Aunque reconoce que otros factores influyen en los accidentes, el análisis se centra principalmente en la velocidad promedio del tráfico y no en variables meteorológicas. De igual manera, el documento no estudia en profundidad las características específicas de las intersecciones, como diferentes tipos de cruces, rotondas o semáforos, ni cómo estos elementos pueden cambiar el riesgo de accidentes. Las intersecciones se mencionan en el contexto de carreteras urbanas, pero no se analizan como una variable independiente dentro del modelo. 

	\item \textbf{Contraste con otra fuente:} Una conclusión particular a la que llegan los autores de Traffic Accidents Analysis on Dry and Wet Road Bends Surfaces in Greater Manchester-UK es que la velocidad excesiva es un factor importante que contribuye a la ocurrencia de accidentes, lo cual coincide con lo planteado en el documento de M. C. Taylor, D. A. Lynam y A. Baruya. Por otra parte, en el texto The Effects of Drivers' Speed on the Frequency of Road Accidents, en la figura 5 se muestra un diagrama con la interrelación de variables que afectan la frecuencia de accidentes. En este se mencionan factores como la iluminación, el clima y otras condiciones de circulación que el estudio no analiza directamente. Al complementarlo con la investigación de Karzan S. Ismael y Noorance A. Razzaq se obtiene un panorama más concreto, ya que este trabajo se enfoca en analizar otros factores del entorno vial. En particular, examina cómo las condiciones de la superficie de la carretera (seca o mojada), la iluminación y el comportamiento del conductor influyen en la ocurrencia de accidentes en curvas. Los resultados muestran que aproximadamente el 41 \% de los accidentes analizados ocurrieron en carreteras mojadas, y que factores como la pérdida de control del vehículo, el deslizamiento y los errores del conductor tienen un papel importante en estos accidentes. Por lo tanto, se puede observar que, aunque ambos documentos tienen enfoques distintos, se complementan y llegan a conclusiones similares sobre la importancia de la velocidad y otros factores en la ocurrencia de accidentes.   
	
	\item \textbf{Evaluación de aplicabilidad (1-5):} La aplicabildad se valora en un 0 de 5 ya que, aunque la teoría es útil para respaldar argumentos de que la velocidad tiene una relación positiva con la frecuencia de accidentes, la metodología en sí del estudio es de carácter comparativo, pues integra resultados de diferentes investigaciones y contextos para establecer dicha relación. De igual forma, difícilmente puede replicarse cuando se carece del dato sobre la velocidad promedio.   
	
	\item \textbf{Pregunta que le haría al autor:} Si su estudio se basa en la relación entre la velocidad promedio del tráfico y la frecuencia de accidentes, ¿cómo podría aplicarse este enfoque en investigaciones que solo cuentan con variables indirectas como el límite de velocidad de la vía o el tipo de carretera? 
	
	\item \textbf{Resumen argumentativo:} El artículo de M. C. Taylor, D. A. Lynam y A. Baruya intenta resolver el problema de cómo entender y cuantificar la relación entre la velocidad del tráfico y la frecuencia de accidentes. Los autores proponen analizar esta relación utilizando la velocidad promedio de los vehículos en la vía lo cual permite observar cómo cambios pequeños en la velocidad general de una vía pueden afectar el número total de accidentes. Sin embargo, su aplicación puede ser limitada en contextos donde no se dispone de datos sobre la velocidad promedio, ya que estudios como el que se pretende realizar solo cuentan con información como límites de velocidad o características de la carretera. A pesar de esto, el estudio proporciona evidencia clara de que incluso pequeñas reducciones en la velocidad promedio pueden generar disminuciones significativas en la frecuencia de accidentes. En el trabajo se utilizará este documento principalmente como base teórica para justificar la importancia de la velocidad como factor que influye en la frecuencia de accidentes, aunque posiblemente no se pueda aplicar exactamente su metodología ya que los datos disponibles no incluyen medidas de velocidad promedio. Además, servirá para entender cómo la velocidad interactúa con otros factores del entorno vial. \\
	\end{itemize} 
	

\textbf{Ficha 3:}

\textbf{Nivel descriptivo} 
\begin{itemize}
    \item \textbf{Título:} Accident risk of road and weather conditions on different road types
    
	\item \textbf{Autores:} Fanny Malin, Ilkka Norros, Saty Innamaa 
    
	\item \textbf{Año:} 2019
    
	\item \textbf{Nombre del tema:} Este estudio pretende evaluar el riesgo de accidentes de tráfico en diferentes condiciones climáticas y tipos de carretera de Finlandia.
	
    \item \textbf{Clasificación:} Metodológica ya que este presenta y aplica un método basado en la distribución de Palm de diferentes condiciones y comparándolo con los resultados de otras distribuciones para analizar que produce mayor riesgo de accidentes    
	
    \item \textbf{Resumen en una oración:} Estudio acerca del riesgo de accidentes de tránsito según las condiciones climáticas y el tipo de carretera. 

    \item \textbf{Argumento central:} Se demuestra que el riesgo de accidentes varía significativamente según las condiciones del clima y el tipo de carretera, siendo mayores el riesgo relativo durante lluvias heladas y cuando la calle se encuentra resbaladiza y además, cuando se trata de autopistas bajo estas condiciones. Se utiliza la probabilidad de Palm para evaluar el riesgo desde la perspectiva del conductor.  
 
    \item \textbf{Problemas con el argumento o el tema:} El periodo de estudio es relativamente muy corto (2 años) y solo basado en los datos de 43 rutas de Finlandia. Además, las categorías de las condiciones del clima se solapan e inclusive algunas de las condiciones son muy extrañas lo que las hace estadísticamente imposible que ocurran.  \\
\end{itemize}

\textbf{Nivel analítico }
\begin{itemize}
	\item \textbf{Conexión con mi proyecto:}Ambos proyectos intentan conectar y encontrar el peso que tienen algunas variables frente el riesgo de causar un accidente de tránsito. Principalmente el clima y las condiciones de la carretera, variables presentes en la base de datos seleccionada. La base de datos a trabajar, describe la carretera, las condiciones de la misma y el clima en el momento del accidente; es decir, si es una carretera de una vía, doble vía, si está seca, mojada, con hielo y si llueve o hace sol, entre otras opciones.  
 
	\item \textbf{Lo que NO dice el autor:} Este estudio toma lugar en Finlandia mientras que el proyecto a trabajar toma datos del Reino Unido, mientras ambos países son Europeos sus condiciones climáticas son muy distintas al igual que sus carreteras. Por ejemplo, el sentido de circulación es distinto, las carreteras en Finlandia están diseñadas para las condiciones invernales severas con gran cantidad de nieve, mientras que el Reino Unido es famoso por la cantidad de lluvia que recibe. Asimismo, el método utilizado de la probabilidad de Palm no es uno que se vaya a tratar posteriormente. Por el lado, este estudio considera el riesgo relativo a ciertos factores a partir de cuánto tiempo ha pasado el conductor en carretera en cierto punto. Esta característica no es brindada en la base de datos seleccionada referente al Reino Unido.  
 
	\item \textbf{Contraste con otra fuente:} El estudio “Accident risk of road and weather conditions on different road types” examina las condiciones ambientales y su impacto relativo sobre el riesgo de un accidente de tránsito donde la variable de velocidad es supuesta dada los tipos de carretera, mientras que en “The effects of drivers' speed on the frequency of road accidents” la velocidad es el factor principal a analizar con respecto a la frecuencia de dichos accidentes y las condiciones del clima son secundarias. Sin embargo, ambos estudios presentan variables como el tipo de carretera y las condiciones de circulación.    
 
	\item \textbf{Evaluación de aplicabilidad:} La aplicabilidad es de un 3, ya que su metodología como tal no es de interés para este proyecto. No se busca realizar una distribución de Palm; sin embargo, la estructura y el análisis realizado en el estudio si es relevante a lo que se desea crear en su proyecto. Además, las conclusiones que realizaron son de suma importancia teórica para los resultados esperados de la investigación. Por ejemplo, como las condiciones nevadas provocan mayor riesgo de accidentes comparados con una lluvia leve o como el riesgo de accidentes aumenta en pista, entre otros.  
 
	\item \textbf{Pregunta que le haría al autor:} ¿Cómo podría adaptarse el análisis relativo de riesgo que ustedes desarrollaron para incluir otras variables como velocidad y además en condiciones climáticas como de infraestructura diferentes (de Finlandia al Reino Unido)?    
 
	\item \textbf{Resumen argumentativo:} Los autores desean resolver el problema de cuantificar y comparar el riesgo de accidentes de tránsito bajo distintas condiciones meteorológicas y tipos de carretera. Esto lo hace desde la perspectiva de cuánto tiempo pasa el conductor en dichas condiciones. Sus resultados acerca de cómo ciertos factores como las condiciones de nieve intensa y las carreteras resbaladizas elevan el riesgo de accidentes son relaciones que se buscarán dentro de este proyecto con las variables respectivas al Reino Unido. Es decir, estas situaciones identificadas en Finlandia donde ciertas condiciones climáticas y la condición de la carretera afectan la frecuencia de accidentes funcionan como base teórica. Sin embargo, la metodología basada en la distribución de Palm y el análisis del tráfico por trayecto no van a ser aplicadas en el contexto del Reino Unido entre 2010-2017. \\  
\end{itemize}


\textbf{Ficha 4:}

\textbf{Nivel descriptivo}
\begin{itemize}
    \item \textbf{Título:} Analysis of Road Traffic Accident using Causation Theory with Traffic Safety Model and Measures.
    
    \item \textbf{Autores:} Dr. Neeta Saxena 
    
    \item \textbf{Año:} 2017 
    
    \item \textbf{Nombre del tema:} Realizan un análisis de las causas de los accidentes en India, donde se utiliza la Teoría del Dominó y la Teoría de Homeostasis del Riesgo para sugerir medidas de seguridad vial.
    
    \item \textbf{Clasificación:} Teórica. El artículo fundamenta su análisis de los accidentes de tráfico con modelos ya establecidos, como la Teoría de Homeostasis de Riesgo y la Teoría del Dominó.
    
    \item \textbf{Resumen en una oración:} El artículo analiza accidentes de tráfico en la India usando Teorías de seguridad vial.
    
    \item \textbf{Argumento central:} El argumento central dice que los accidentes de tráfico no son eventos solamente accidentales, sino el resultado de una cadena de factores y comportamientos humanos regulados por la percepción del riesgo, los cuales pueden ser prevenidos mediante intervenciones estructurales, educativas y de comportamiento. 
    
    \item \textbf{Problemas con el argumento o el tema:} La autora señala que los accidentes son fenómenos aleatorios difíciles de predecir con exactitud mediante modelos. Además, menciona que actualmente falta información esencial para probar las ecuaciones de la homeostasis del riesgo y que los datos sobre accidentes que solo involucran daños materiales o lesiones no mortales tienden no ser reportados de manera confiable.  \\
\end{itemize}

\textbf{Nivel analítico}

\begin{itemize}
    \item \textbf{Conexión con mi proyecto:}El artículo es relevante para el proyecto, pues analiza que los accidentes no se dan por culpa de un solo factor, sino de una secuencia de factores que al final causan el accidente, lo cual ayuda a justificar el por qué analizar varios factores que pueden considerarse causantes de accidentes de tráfico.
    
	\item \textbf{Lo que NO dice el autor:} En el artículo el trabajo está centralizado en India, por lo que las conductas de las personas, la cultura, el tipo de carreteras, leyes, entre otros diferentes factores, no van a ser los mismos en ambos países, entonces el análisis y recomendaciones hechas no se pueden aplicar de igual forma.
    
	\item \textbf{Contraste con otra fuente:} El artículo “Accident risk of road and weather conditions on different road types” se logra conectar de forma crítica con el de Saxena (2017). Mientras Saxena se enfoca en la Teoría de la Homeostasis del Riesgo, el comportamiento humano ajustándose a la seguridad, Malin et al. proporcionan una observación empírica de cómo el entorno físico, en este caso clima y tipo de vía,  altera el riesgo de causar accidentes de tráfico. Un punto de contradicción es que Saxena sugiere que las mejoras en la vía podrían ser anuladas por conductores más imprudentes, mientras que Malin demuestra que, independientemente del comportamiento, ciertas infraestructuras, tales como las autopistas bajo lluvia helada, presentan riesgos que superan la capacidad de adaptación del conductor. Entonces un artículo logra resolver el vacío del otro, pues Saxena explica el "porqué" humano y Malin el "dónde y cuándo" desde la parte ambiental.
    
	\item \textbf{Evaluación de aplicabilidad (1-5):} 4/5. La metodología de la Teoría del Dominó es bastante aplicable, pues da una herramienta que es capaz de categorizar los factores que brinda la base de datos de forma lógica. En el análisis se pueden categorizar en el ambiente social, fallas de la persona conductora, el acto inseguro y la severidad del accidente.
    
	\item \textbf{Pregunta que le haría al autor:} Considerando que la base de datos refleja ,como usted lo menciona, un sistema triple E (Ingeniería, Educación y Cumplimiento)  más desarrollado, ¿cómo ajustaría la ecuación accidentabilidad para aislar los efectos que no son tan prominentes dentro de Reino Unido?
    
	\item \textbf{Resumen argumentativo:} La autora intenta resolver el problema de la siniestralidad vial en la India usando varias teorías y modelos para poder proponer soluciones para que estos disminuya. Además,  se analiza cómo una secuencia de fallos mecánicos y humanos logra desencadenar un accidente inevitable una vez que esta cadena comienza. La solución propuesta a través de la Teoría del Dominó funciona para el contexto en el Reino Unido, ya que permite organizar múltiples variables en una jerarquía de causas directas e indirectas. El elemento que se puede usar es la interrupción de la secuencia causal, utilizando los datos de 2010 a 2017 para identificar cuál es la "ficha" más recurrente en los accidentes, por ejemplo, si es la falla humana o las condiciones ambientales. También, basándose en la Teoría del Dominó, se fundamentará las propuestas de prevención argumentando que la intervención en esa “ficha” específica es la forma más efectiva de detener la frecuencia de los siniestros.
	\end{itemize}
	
\textbf{Ficha 5:}

\textbf{Nivel descriptivo}
\begin{itemize}
    \item \textbf{Título:} Análisis estadístico de tablas de contingencia y chi-cuadrado para medir el flujo migratorio de
    origen y destino en el Ecuador año 2018.
    
    \item \textbf{Autores:}  David Lastre Baquerizo, Melanie Páez Santana, y Olga López Tumbaco.
    
    \item \textbf{Año:} 2019
    
    \item \textbf{Nombre del tema:} Aplicación de tablas de contingencia y prueba chi-cuadrado para analizar el flujo 
    migratorio en Ecuador.
    
    \item \textbf{Clasificación:} Metodológico. El artículo explica detalladamente el funcionamiento de las tablas de contingencia y de 
    la prueba de independencia Chi-Cuadrado.
    
    \item \textbf{Resumen en una oración:} Uso de pruebas de independencia Chi-cuadrado para analizar patrones de flujo migratorio en Ecuador.
    
    \item \textbf{Argumento central:} El artículo sostiene que el uso de herramientas estadísticas como las tablas de contingencia y la prueba 
    chi-cuadrado permite analizar de manera objetiva y eficiente los patrones de migración entre origen y destino, facilitando la interpretación 
    de relaciones entre variables y apoyando la toma de decisiones basada en datos.
    
    \item \textbf{Problemas con el argumento o el tema:} El estudio muestra limitaciones relacionadas con la disponibilidad y calidad de los datos 
    históricos, además de la dependencia de supuestos estadísticos que pueden afectar la precisión del análisis. También puede existir dificultad para 
    explicar la naturaleza social del fenómeno migratorio únicamente mediante métodos cuantitativos.\\
\end{itemize}


\textbf{Nivel analítico}

\begin{itemize}
    \item \textbf{Conexión con mi proyecto:} El artículo se conecta con el presemte proyecto en la metodología que utiliza,
    estas son tablas de contingencia y pruebas de Chi-Cuadrado para encontrar patrones de migración (de salida o entrada) 
    por sexo, nacionalidad, motivo de viaje, entre otras. 
    
	\item \textbf{Lo que NO dice el autor:} El artículo no responde a ningún aspecto de la pregunta de investigación, ni permite identificar relaciones 
    específicas entre variables propias de este fenómeno. Su principal limitación para este proyecto es que, aunque desarrolla herramientas metodológicas como 
    las tablas de contingencia y la prueba chi-cuadrado, no las aplica en el contexto de accidentes de tránsito. 
    
	\item \textbf{Contraste con otra fuente:} El artículo de Lastre et al. propone el uso de pruebas de independencia Chi-Cuadrado para analizar la relación 
    entre variables categóricas en un contexto de migración. Por otro lado, el artículo de \textcite{MONTANA2025} se enfoca en la prueba de homogeneidad 
    Chi-Cuadrado, en el ámbito de estudios clínicos. Ambos documentos se complementan ya que abordan distintas aplicaciones de la misma familia de métodos 
    estadísticos. El primero se centra en encontrar una relación significativa entre variables, es decir, cuantificar la dependencia o independencia de las 
    variables, mientras que el segundo se enfoca en determinar si existen diferencias o similitudes en la distribución de distintas categorías de las variables. 
    Esto es relevante para esta investigación pues no solo resulta relevante encontrar relación entre distintos factores y la frecuencia de accidentes de tránsito, sino 
    también comparar cómo varían y se distribuyen los accidentes entre distintos grupos y condiciones. Además, el segundo artículo ayuda a resolver una limitación 
    del primero, ya que amplía el enfoque metodológico al introducir la prueba de homogeneidad como herramienta adicional para el análisis de variables categóricas. 
    Sin embargo, ambos comparten el problema de no abordar directamente el tema de accidentes de tránsito, por lo que su aporte se mantiene en el plano metodológico.
    
	\item \textbf{Evaluación de aplicabilidad (1-5):} La metodología es aplicable porque tanto el artículo como el proyecto trabajan con variables categóricas, lo que 
    hace pertinente el uso de tablas de contingencia y la prueba chi-cuadrado.  Sin embargo, no se asigna un 5 porque el artículo no aborda las particularidades de los 
    datos de accidentes de tránsito. 
    
	\item \textbf{Pregunta que le haría al autor:} ¿Cómo recomienda adaptar la prueba chi-cuadrado cuando se trabaja con datos de accidentes de tránsito que pueden presentar 
    baja frecuencia en ciertas categorías o posible dependencia entre observaciones?
    
	\item \textbf{Resumen argumentativo:} El autor busca resolver el problema de cómo analizar la relación entre variables categóricas mediante herramientas estadísticas 
    como las tablas de contingencia y la prueba chi-cuadrado para explicar el flujo migratorio en Ecuador. Su propuesta funciona porque estas técnicas permiten cuantificar 
    la asociación entre variables, permitiendo interpretar patrones dentro de los datos. En este proyecto, la solución es adecuada en cierta medida, ya que las variables 
    con las que se trabajarán también son categóricas y el objetivo es identificar relaciones entre factores ambientales y la frecuencia de accidentes. Sin embargo, el enfoque 
    del artículo no considera la complejidad propia de los accidentes de tránsito, como su naturaleza multicausal o la presencia de dependencias entre observaciones. A pesar de 
    esto, la metodología aplicada en el artículo sigue siendo aceptable como herramienta de análisis. Específicamente, el uso de tablas de contingencia permitirá organizar y 
    visualizar las frecuencias cruzadas entre variables importantes. Asimismo, la prueba chi-cuadrado será utilizada para determinar si existe evidencia estadística de asociación 
    entre estas variables. Por lo tanto, el principal aporte del artículo al proyecto es proporcionar una base metodológica completa para el análisis de las variables consideradas.\\
\end{itemize} 

\textbf{Ficha 6:}

\textbf{Nivel descriptivo}
\begin{itemize}
    \item \textbf{Título:} UNA INTRODUCCIÓN A LA IMPUTACIÓN DE VALORES PERDIDOS
    
    \item \textbf{Autores:} Lelly Useche y Dulce Mesa
    
    \item \textbf{Año:} 2006 
    
    \item \textbf{Nombre del tema:} Uso de diferentes métodos estadísticos y 
    matemáticos para lograr imputar de forma correcta los valores perdidos de 
    una base de datos, así como explicar sus causas y tipos de errores.
    
 \item \textbf{Clasificación:} Metodológica. El artículo explica varias herramientas
 matemáticas y estadísticas para poder encontrar la mejor manera de imputar los
 valores perdidos, además de describir las ventajas y los criterios para poder 
 seleccionar una técnica de imputación específica

    
    \item \textbf{Resumen en una oración:} El artículo analiza los
    valores perdidos y describe métodos para imputarlos.
    
    \item \textbf{Argumento central:} El documento busca que se
    comprenda que la ausencia de los datos genera sesgos y pérdida 
    de eficiencia en los análisis., por lo que la imputación busca 
    rellenar esos vacíos con valores que tengan sentido para recuperar 
    la estructura de la base de datos y así usar los métodos tradicionales 
    estadísticos para resolver el problema de interés.
    
    \item \textbf{Problemas con el argumento o el tema:} Los autores 
    dicen que ninguna técnica de imputación es perfecta, puesto que 
    pueden introducir sesgos en las estimaciones, alterar la relación 
    entre variables o cambiar la varianza de los datos.  \\
\end{itemize}

\textbf{Nivel analítico}

\begin{itemize}
    \item \textbf{Conexión con mi proyecto:} El documento es de 
    mucha utilidad pues ayuda a poder realizar la limpieza de datos 
    de la base de nuestro proyecto tomando en cuenta cuales variables es mejor o no 
    imputar y saber las consecuencias tanto positivas como negativas del método de imputación que se vaya a escoger para los valores perdidos.
    
	\item \textbf{Lo que NO dice el autor:} Los autores no explican cómo tratar estos 
	valores perdidos con variables de carácter geográfico o temporales. El documento no tiene 
	información de qué hacer con datos donde el error tal vez no sea por recolección, sino 
	información no disponible, que por sí misma podría ser un factor de riesgo.    
	\item \textbf{Contraste con otra fuente:} Este documento se complementa con el 
	estudio de \textcite{KARZAN2017} acerca accidentes en Manchester. Mientras que Karzan 
	utiliza la prueba Chi-Cuadrado para encontrar relaciones entre el clima y los choques, 
	Useche explica que esos cálculos solo serán confiables si primero se resuelve el 
	problema de los datos faltantes. Es decir, el primer texto da la herramienta para 
	limpiar la base de datos, asegurando que los valores perdidos no distorsionen el 
	análisis, y el segundo da la técnica estadística para analizar los datos una vez 
	que estén completos. Uno resuelve el problema técnico de la base de datos que el otro da por sentado.
    
	\item \textbf{Evaluación de aplicabilidad (1-5):} 5/5. La aplicabilidad de estas 
	metodologías es perfecta, pues en la base de datos existen valores perdidos y la 
	forma en la que estos se manejen va a afectar de una u otra manera el análisis 
	estadístico que se hará. Además, el documento advierte que el ignorar totalmente 
	los valores perdidos puede invalidad el análisis de variables multivariantes. 
	Por otro lado, el objetivo del proyecto es identificar situaciones de riesgo, 
	por lo que la imputación que se realice afectará directamente en los resultados 
	que darán los análisis para poder dar las recomendaciones pertinentes.
    
	\item \textbf{Pregunta que le haría al autor:} En el caso de bases de datos 
	masivas como la base de datos del proyecto, donde mucha información es descriptiva 
	(como el clima, el tipo de carretera o qué tan grave fue el choque), ¿es mejor 
	usar métodos que predicen el dato faltante basándose en las otras respuestas, o
	es más seguro usar promedios generales para no alterar los resultados finales?
    
	\item \textbf{Resumen argumentativo:} Los autores buscan solucionar el problema 
	de los datos faltantes en las bases de datos, ya que estos pueden dar resultados 
	falsos o incompletos. Esto es importante para el proyecto pues, en los accidentes 
	de tránsito del Reino Unido, la información suele venir incompleta debido a la 
	confusión que hay en el lugar del choque. La técnica de llenar estos espacios con 
	valores lógicos es de gran utilidad, porque así no se tiene que descartar accidentes 
	enteros solo por la falta de un dato. Del documento, se usará la guía para poder 
	imputar aquellos valores faltantes de aquellas variables que se considere importante 
	de modificar. Esto para lograr que las propuestas para prevenir accidentes sean lo más 
	confiables posibles y basadas en información que tenga concordancia entre sí.\\
	\end{itemize}

\textbf{Ficha 7:}

\textbf{Nivel descriptivo}
\begin{itemize}
    \item \textbf{Título:} Prueba de Chi cuadrado de homogeneidad en estudios clínicos: una 
    herramienta para analizar diferencias entre tratamientos
    
    \item \textbf{Autores:} R.M. Montaña, Á. Roco-Videla, A.R. Nieves y S.V. Flores.
    
    \item \textbf{Año:} 2024 
    
    \item \textbf{Nombre del tema:} El artículo explica el uso de la prueba de Chi-cuadrado 
    de homogeneidad para determinar si existen diferencias significativas en las distribuciones 
    de frecuencia entre distintos grupos.
    
 \item \textbf{Clasificación:} Metodológica. El artículo funciona como una guía de 
 procedimientos que detalla los pasos técnicos, desde el planteamiento de hipótesis 
 hasta el cálculo de residuos estandarizados en software estadístico (SPSS).

    
    \item \textbf{Resumen en una oración:} El documento describe el procedimiento 
    técnico para aplicar pruebas de homogeneidad y comparar grupos independientes.
    
    \item \textbf{Argumento central:} La prueba de Chi-cuadrado de homogeneidad es 
    importante para verificar si diferentes poblaciones o grupos (en este caso, 
    tratamientos médicos) se comportan de manera similar respecto a una variable 
    cualitativa, lo que permite tomar decisiones basadas en la estadística.
    
    \item \textbf{Problemas con el argumento o el tema:} Los autores mencionan que para que 
    la prueba sea válida depende de que las muestras sean independientes y que los valores 
    esperados en las celdas de la tabla de contingencia sean superiores a 5.\\
\end{itemize}

\textbf{Nivel analítico}

\begin{itemize}
    \item \textbf{Conexión con mi proyecto:} En el análisis de la base de datos, esta
    prueba es importante para comparar si la severidad de los accidentes se distribuye 
    de forma homogénea en diferentes regiones o bajo distintos climas. Por ejemplo, 
    se puede usar esta técnica para probar si la proporción de accidentes fatales es 
    la misma en carreteras secas que en carreteras húmedas, tal como los autores 
    comparan la eficacia de dos analgésicos.    
	\item \textbf{Lo que NO dice el autor:} El documento se enfoca en muestras 
	pequeñas de estudios clínicos (215 casos). En el caso de la base de datos, 
	se trabajará con una gran cantidad de datos, donde casi cualquier diferencia 
	mínima resulta estadísticamente significativa. Los autores no discuten el tamaño 
	del efecto o la importancia práctica frente a la significancia estadística con una base de datos masiva.
	\item \textbf{Contraste con otra fuente:} El estudio se relaciona con el
	estudio de \textcite{USECHE2006}, pues Useche nos habla sobre como 
	imputar los datos para que la base de datos esté lo más completa posible, tomando 
	en cuenta diferentes factores, mientras que \textcite{MONTANA2025} nos da un manual 
	técnico para procesar los datos ya corregidos usando la prueba de Chi-Cuadrado. 
	Entonces \textcite{USECHE2006} garantiza que la información no esté sesgada para
	con \textcite{MONTANA2025} se pueda llegar conclusiones significativas de la información extraída.
    
	\item \textbf{Evaluación de aplicabilidad (1-5):} 4/5.  El método es bastante 
	aplicable para la base de datos en la que se trabajará, pues se puede usar para 
	variables categóricas como el clima e iluminación. Además de que da un paso a 
	paso de como se puede analizar ese tipo de datos. Una desventaja es que el enfoque 
	de ambos estudios es muy diferente y no se centra en los datos masivos, por lo que 
	la eficacia de este no se puede verificar al máximo.
    
	\item \textbf{Pregunta que le haría al autor:} ¿Cómo se vería afectada la potencia de 
	esta prueba de homogeneidad si en lugar de comparar solo dos grupos (como los analgésicos), 
	comparamos diez tipos diferentes de superficies viales en una tabla de contingencia mucho más grande?    
	\item \textbf{Resumen argumentativo:} Los autores buscan facilitar la comparación 
	de grupos en investigaciones mediante una herramienta que detecta si las diferencias 
	observadas son reales o producto del azar. Esta solución es perfecta para el proyecto 
	porque permite segmentar la base de datos y verificar si el riesgo de accidente cambia 
	según el entorno. La metodología funciona para el contexto, ya que el análisis de residuos 
	estandarizados corregidos que se menciona permite identificar exactamente qué combinación es 
	la que dispara el riesgo, por ejemplo, lluvias e intersecciones. Del documento, se usará la 
	guía para interpretar los resultadas y lograr asegurar que las propuestas de prevención de 
	accidentes tengan el rigor estadístico que se exige.\\
	\end{itemize}


\textbf{Ficha 8:}

\textbf{Nivel descriptivo}
\begin{itemize}
    \item \textbf{Título:} Major theories of construction accident causation models: a literature review 

    \item \textbf{Autores:}  Seyyed Shahab, Zahra Jabbarani Torghabeh
    
    \item \textbf{Año:} 2012
    
    \item \textbf{Nombre del tema:} Las teorías principales sobre la causalidad de accidentes y sus análisis 
    
 \item \textbf{Clasificación:} Teórica, porque el documento revisa y analiza 
 diversas teorías y modelos de causalidad de accidentes para explicar cómo y por qué ocurren. 
    
    \item \textbf{Resumen en una oración:} Análisis de las teorías que 
    explican las causas de los accidentes de construcción y su prevención
    
    \item \textbf{Argumento central:} El documento argumenta que 
    aunque existen múltiples teorías de causalidad de accidentes donde se enfatizan los factores 
    humanos, condiciones inseguras y gestión, estas teorías no generan estrategias para reducir 
    el riesgo en obras de construcción, por lo que se propone que los sistemas de gestión de
    seguridad incorporen las imperfecciones naturales en sus directrices para prevenir accidentes. 
    
    \item \textbf{Problemas con el argumento o el tema:} Un problema presentado es que las 
    teorías existentes presentan una visión simplificada que atribuye la mayoría de los accidentes a 
    errores humanos sin considerar que no todos los riesgos pueden eliminarse con un cambio de comportamiento. 
    Además, otro problema es que estas teorías no brindan directrices prácticas 
    que se puedan aplicar y que todo accidente es evitable. 

\end{itemize}

\textbf{Nivel analítico}

\begin{itemize}
    \item \textbf{Conexión con mi proyecto:}  Este documento se quiere utilizar 
    como parte de las teorías utilizadas, se busca tomar la parte conceptual que 
    brinda de las distintas teorías acerca de las causas de los accidentes como 
    el Modelo de Múltiples Causas de Petersen, La Teoría de Dominio actualizada 
    de Weaver, etc. Para posteriormente tomar la información conceptual sobre 
    cada una de ellas y relacionarla con accidentes de tránsito.
    
	\item \textbf{Lo que NO dice el autor:}  Dado que este documento habla de accidentes de 
	construcción, los elementos que contextualizan en sus modelos son distintos a 
	los necesarios si se quiere analizar accidentes de tránsito, como se desea. 
	Por lo que dadas las explicaciones tanto generales como específicas se puede 
	relacionar con los accidentes de tránsito, por ejemplo en el modelo de Petersen 
	se consideran las reglas impuestas en el lugar para los accidentes de construcción 
	se puede considerar el manejo de maquinaria pesada, en nuestro caso estas reglas 
	serían el límite de velocidad de la vía, si se puede adelantar en el carril, etc. 
	
	\item \textbf{Contraste con otra fuente:}  Podemos realizar una comparación con el 
	artículo Analysis of Road Traffic Accident using Causation Theory with Traffic 
	Safety Model and Measures de Saxena. Ambos documentos coinciden en que los accidentes 
	no son eventos aislados y que resultan de una secuencia o interacción de múltiples 
	factores. Además, enfatizan que tanto el comportamiento humano como el ambiente 
	interactúan de tal manera que ciertos elementos aumentan el riesgo de que suceda 
	un accidente. Sexena señala la dificultad de predecir accidentes debido a su aleatoriedad 
	y la falta de datos confiables, mientras que el artículo de Shahab et al, menciona que 
	la dificultad principal es que los modelos de causación de accidentes se enfocan demasiado 
	en el factor humano. Un inconveniente que presentan ambos documentos es su aplicabilidad
	a nuestro proyecto, ya que aunque el trabajo de Saxena si se basa en accidentes de 
	tránsito estos son en la India por lo tanto las condiciones, regulaciones y costumbres 
	son algunas distintas a aquellas del Reino Unido y, en el caso del realizado por Shahab 
	los accidentes trabajados son de construcción. Por lo tanto, ninguno de ellos es completamente aplicable a nuestro trabajo. 
	
	\item \textbf{Evaluación de aplicabilidad (1-5):} Su aplicabilidad es de 3, dado 
	que la teoría es la única parte de utilidad. Sin embargo, toda la discusión 
	donde aplican las distintas teorías de cómo se producen accidentes son aplicadas 
	a un tema que no es de nuestro interés (accidentes de construcción), por lo 
	que solo podemos usarlo como inspiración.  
	
	\item \textbf{Pregunta que le haría al autor:}  Si cambiara el enfoque de su 
	proyecto a accidentes automovilísticos, ¿qué elementos tomaría en cuenta? 
	
	\item \textbf{Resumen argumentativo:} El problema central que se intenta 
	resolver es la falta de comprensión con respecto a las causas de los accidentes
	de construcción, donde muchas de las teorías existentes atribuyen de forma 
	excesiva los accidentes a errores humanos sin considerar que hay riesgos 
	inherentes imposibles de eliminar por completo. Se propone desarrollar 
	modelos y sistemas de gestión de seguridad que incorporen tanto las 
	imperfecciones humanas como las amenazas presentes en un entorno de construcción 
	para prevenir accidentes de una manera más efectiva. Sin embargo, este artículo 
	se centra en prevenir accidentes de construcción lo que limita su aplicabilidad
	a otros tipos de accidentes debido a las diferencias en los factores ambientales, 
	operacionales y de regulación. Para este proyecto acerca de accidentes de tránsito, 
	se utilizaran los conceptos como el Modelo de Múltiples Causas de Petersen y 
	la Teoría de Dominó actualizada de Weaver. A partir de la base teórica se adaptará 
	para incluir factores relevantes de accidentes de tránsito. 
	\end{itemize}
	
	
	
	
	
	
	
	
\subsection{Análisis de modelación}


\subsection{Propuesta metodológica}




\subsection{Construcción de fichas de resultados}
\textbf{Ficha 1:}

\textbf{Nivel descriptivo}
\begin{itemize}
    \item \textbf{Títular:} 
    
    \item \textbf{Nombre del hallazgo/resultado:} 
    
    \item \textbf{Resumen en una oración:}  
    
    \item \textbf{Método o análisis que lo produjo:} 
    
    \item \textbf{Evidencia:} 
    

\end{itemize}

\textbf{Nivel analítico}

\begin{itemize}
    \item \textbf{Conexión con la pregunta de investigación:} 
    
	\item \textbf{Constraste con la literatura:} 
	
	\item \textbf{Lo que NO explica este resultado:}   
    
	\item \textbf{Implicación para el siguiente paso:} 
	
	
\end{itemize} 










\subsection{Ordenamiento de los elementos del reporte}
\subsection{La imagen que cuenta la historia}


\section{Arco narrativo del proyecto}
\section{Parte de escritura}
\section{Parte de reflexión}
\section{Diario de aprendizaje}





